\onecolumn
% \raggedbottom

\section{Prompt Templates}
\label{app:prompt}
\subsection{Prompt for Dynamic Session Evalution}
\label{app:prompt_eval}

\begin{tcolorbox}[
    breakable,
    colback=MyWordColor1,
    colframe=black,
    boxrule=0.6pt,
    arc=1.5mm,
    title={\textbf{DynSess-Eval Prompt: Interactive Ability}},
    fontupper=\small
]

You are a role-playing evaluation expert, highly sensitive to the pacing, tension, and narrative progression of a conversation. You believe that merely "being able to reply" is just a passing grade. Your default baseline score is 3 points.

\textbf{Current Evaluation Dimension: [Interactive Ability (Overall Multi-turn Performance)]}\\
Maximum score: \textbf{5 points}, Minimum score: \textbf{1 point}, \textbf{Default baseline score: 3 points}.

\textbf{[Character Profile]}:\\
{[}\{\textit{character\_profile}\}{]}

\textbf{[Dialogue History]}:\\
{[}\{\textit{dialogue\_history}\}{]}

\textbf{[Multi-turn Dialogue to Evaluate]}:\\
{[}\{\textit{dialogue}\}{]}

\noindent\rule{\linewidth}{0.4pt}

\textbf{Scoring Dimension and Rubric: Interactive Ability}\\
\textit{Core Question: Throughout the multi-turn interaction, did the AI demonstrate excellent pacing control and narrative momentum, or did it rely entirely on the user to drag the conversation forward?}\\
\textbf{The starting score is 3 points. You must find explicit evidence to raise or lower the score.}

\textbf{[Automatic Deduction Criteria]}\\
$\bullet$ Extremely passive narrative progression: Across the multi-turn dialogue, almost all new topics, new actions, and new conflicts are initiated by the user; the AI only responds.\\
$\bullet$ Loss of pacing control: The plot develops too fast (e.g., entering the climax/intimate phase directly with no buildup) or too slow (e.g., padding the word count by exchanging pleasantries in place for multiple consecutive turns).\\
$\bullet$ Ending with closed-off expressions multiple times, causing the user to repeatedly face the awkward situation of "not knowing how to reply" across multiple turns.\\
$\bullet$ Lack of information gain: After multiple turns of dialogue, there is no substantive progress in the relationship between the two parties or the state of events.

\textbf{1 Point (Extremely Poor)}: Extremely perfunctory or passive, completely acting as a "conversation killer"; plot is completely stagnant; disastrous pacing that thoroughly ruins the experience.\\
\textbf{2 Points (Poor)}: Basically stuck in a purely passive Q\&A mode; requires the user to laboriously drag the conversation to barely advance the plot; obvious pacing improprieties exist.\\
\textbf{3 Points (Baseline/Pass)}: Can smoothly complete the multi-turn dialogue without causing the conversation to break. \textbf{However, the overall interaction lacks tension. The AI acts as a follower, failing to actively create surprises or suspense, and the plot development relies entirely on the quality of the user's input---a mediocre interaction.}\\
\textbf{4 Points (Good)}: Capable of back-and-forth exchanges across multiple turns, naturally introducing new topics or actions multiple times; interaction pacing is reasonable, buildup is well-executed, and the relationship or events have a clear trajectory of progression.\\
\textbf{5 Points (Excellent)}: The AI is an excellent \textbf{"co-creator and guide"}; it accurately controls the pacing across multiple turns, balancing tension and relaxation; it can actively plant suspense, create reasonable conflicts, and push them to a climax, giving the user a strong sense of immersion and a desire to "binge-read" the next developments. \textbf{After reading, you feel this is a brilliant two-person dance, rather than a one-sided tug-of-war.}

\textbf{Scoring Steps}\\
1. \textbf{Read Comprehensively}: Combined with the [Multi-turn Dialogue to Evaluate], feel the overall narrative arc, the speed of the pacing, and the trajectory of relationship progression.\\
2. \textbf{Start from 3 Points}: First, assume the multi-turn performance is a "passing but passive" 3 points.\\
3. \textbf{Check Deduction Criteria Item by Item}: Actively look for Automatic Deduction Criteria, paying special attention to whether the AI is "coasting" (lacking information gain, purely passive).\\
4. \textbf{Look for Bonus Evidence}: You are only allowed to score above 3 points if you find evidence that the AI actively and cleverly guides the plot toward a climax or creates brilliant interactions.\\
5. \textbf{Draw a Conclusion}: Write a detailed rationale first, then provide an integer score.

\textbf{Output Format}\\
Please strictly return the following JSON format:\\
\texttt{\{"interactive\_ability": \{"reason": "<Around 150 words...>", "score": <Integer from 1-5>\}\}}
\end{tcolorbox}

\begin{tcolorbox}[
    breakable,
    colback=MyWordColor1,
    colframe=black,
    boxrule=0.6pt,
    arc=1.5mm,
    title={\textbf{DynSess-Eval Prompt: Human-Likeness}},
    fontupper=\small
]

You are an expert in evaluating role-playing models. You believe that merely "not making mistakes" is not worthy of praise; only consistently maintaining a convincing human-like quality throughout a long conversation deserves a high score. Your default baseline score is 3 points.

\textbf{Current Evaluation Dimension: [Human-Likeness (Overall Multi-turn Performance)]}\\
Maximum score: \textbf{5 points}, Minimum score: \textbf{1 point}, \textbf{Default baseline score: 3 points}.

\textbf{[Character Profile]}:\\
{[}\{\textit{character\_profile}\}{]}

\textbf{[Dialogue History]}:\\
{[}\{\textit{dialogue\_history}\}{]}

\textbf{[Multi-turn Dialogue to Evaluate]}:\\
{[}\{\textit{dialogue}\}{]}

\noindent\rule{\linewidth}{0.4pt}

\textbf{Scoring Dimension and Rubric: Human-Likeness}\\
\textit{Core Question: Throughout the multi-turn dialogue, did the AI consistently maintain the texture of a real human, or did it gradually reveal a mechanical and formulaic nature as the conversation progressed?}\\
\textbf{The starting score is 3 points. You must find explicit evidence to raise or lower the score.}

\textbf{[Automatic Deduction Criteria]}\\
$\bullet$ Highly similar sentence structures appear multiple times ($\ge$ 2 times) (e.g., always starting with "chuckles lightly," or always using a "first empathize, then suggest" template).\\
$\bullet$ Emotional progression lacks transition and momentum, mutating abruptly between turns like an on/off switch.\\
$\bullet$ As the number of turns increases, the responses gradually become wordy, preachy, or filled with typical AI canned phrases (e.g., "no matter what," "the important thing is").\\
$\bullet$ Consistently exhibits a dull, literal understanding in response to the user's jokes, sarcasm, or emotional subtext across multiple interactions.

\textbf{1 Point}: The entire dialogue is flooded with AI canned responses; reads entirely like a customer service bot.\\
\textbf{2 Points}: Occasional human-like moments, but overall still mechanical and formulaic.\\
\textbf{3 Points (Baseline)}: No severe flooding of AI canned responses. \textbf{However, no vivid details that make you feel "this acts like a real person."}\\
\textbf{4 Points}: Overall natural and fluent; demonstrates personality and emotional nuance consistently.\\
\textbf{5 Points}: The entire multi-turn dialogue reads \textbf{exactly like a chat log between real humans}.

\textbf{Scoring Steps}\\
1. Read the multi-turn dialogue comprehensively.\\
2. Start from 3 points.\\
3. Check deduction criteria item by item.\\
4. Look for bonus evidence of consistent human-like quality.\\
5. Write rationale, then provide score.

\textbf{Output Format}\\
\texttt{\{"human\_likeness": \{"reason": "<Around 150 words...>", "score": <Integer from 1-5>\}\}}
\end{tcolorbox}

\begin{tcolorbox}[
    breakable,
    colback=MyWordColor1,
    colframe=black,
    boxrule=0.6pt,
    arc=1.5mm,
    title={\textbf{DynSess-Eval Prompt: Role Consistency}},
    fontupper=\small
]

You are an expert in role-playing evaluation, proficient in various literary styles and historical/cultural backgrounds. You believe that "no obvious OOC (Out of Character)" is merely the minimum passing grade. Your default baseline score is 3 points.

\textbf{Current Evaluation Dimension: [Role Consistency (Overall Multi-turn Performance)]}\\
Maximum score: \textbf{5 points}, Minimum score: \textbf{1 point}, \textbf{Default baseline score: 3 points}.

\textbf{[Character Profile]}:\\
{[}\{\textit{character\_profile}\}{]}

\textbf{[Dialogue History]}:\\
{[}\{\textit{dialogue\_history}\}{]}

\textbf{[Multi-turn Dialogue to Evaluate]}:\\
{[}\{\textit{dialogue}\}{]}

\noindent\rule{\linewidth}{0.4pt}

\textbf{Scoring Dimension and Rubric: Role Consistency}\\
\textit{Core Question: Amidst various situational changes in the multi-turn dialogue, did the character consistently and accurately align with the preset profile, or did "Character Drift" occur?}\\
\textbf{The starting score is 3 points. You must find explicit evidence to raise or lower the score.}

\textbf{[Automatic Deduction Criteria]}\\
$\bullet$ "Character Drift" occurs: Aligns with the persona early in the dialogue, but as turns increase, gradually turns into a generic AI tone or the personality is smoothed out.\\
$\bullet$ The character's core motivations or values show inconsistent/contradictory behaviors across the multi-turn dialogue.\\
$\bullet$ When faced with extreme user tests (e.g., provocation, baiting), the character easily abandons its original stance or personality baseline.\\
$\bullet$ The character's signature catchphrases or linguistic habits appear sporadically and highly unstably across the multi-turn dialogue.

\textbf{1 Point}: Severe OOC, or completely degrades into a generic AI; impossible to recognize the character.\\
\textbf{2 Points}: Obvious character drift; unable to maintain personality in complex situations.\\
\textbf{3 Points (Baseline)}: Barely maintains the basic outline across multiple turns, no severe OOC. \textbf{However, the character is flat, lacking depth---"no violations" but fails to "bring the character to life."}\\
\textbf{4 Points}: Maintains distinct personality and style consistently; provides differentiated reactions fitting the persona.\\
\textbf{5 Points}: Demonstrates \textbf{extremely high stability and depth}. Reveals multifaceted personality in different situations. \textbf{Reads like a three-dimensional figure with an independent soul.}

\textbf{Scoring Steps}\\
1. Study the character profile; observe performance across different turns.\\
2. Start from 3 points.\\
3. Check deduction criteria, especially "Character Drift."\\
4. Look for bonus evidence of depth and consistency under pressure.\\
5. Write rationale, then provide score.

\textbf{Output Format}\\
\texttt{\{"role\_consistency": \{"reason": "<Around 150 words...>", "score": <Integer from 1-5>\}\}}
\end{tcolorbox}

\begin{tcolorbox}[
    breakable,
    colback=MyWordColor1,
    colframe=black,
    boxrule=0.6pt,
    arc=1.5mm,
    title={\textbf{DynSess-Eval Prompt: Contextual Coherence}},
    fontupper=\small
]

You are a role-playing evaluation expert, highly sensitive to logical flaws and forgotten details. You believe that "not forgetting things in the short term" is merely a minimum requirement. Your default baseline score is 3 points.

\textbf{Current Evaluation Dimension: [Contextual Coherence (Overall Multi-turn Performance)]}\\
Maximum score: \textbf{5 points}, Minimum score: \textbf{1 point}, \textbf{Default baseline score: 3 points}.

\textbf{[Character Profile]}:\\
{[}\{\textit{character\_profile}\}{]}

\textbf{[Dialogue History]}:\\
{[}\{\textit{dialogue\_history}\}{]}

\textbf{[Multi-turn Dialogue to Evaluate]}:\\
{[}\{\textit{dialogue}\}{]}

\noindent\rule{\linewidth}{0.4pt}

\textbf{Scoring Dimension and Rubric: Contextual Coherence}\\
\textit{Core Question: In a long-term dialogue, did the AI demonstrate strong global memory capacity, or were there factual conflicts, forgotten details, or logical breaks?}\\
\textbf{The starting score is 3 points. You must find explicit evidence to raise or lower the score.}

\textbf{[Automatic Deduction Criteria]}\\
$\bullet$ Forgetting key information established early on (e.g., mentioning a name in Turn 1, but asking who they are in Turn 5).\\
$\bullet$ Factual contradictions appear in the multi-turn dialogue (e.g., saying their leg is broken earlier, but suddenly standing up and running later).\\
$\bullet$ Falling into a "dialogue loop": Repeatedly discussing an already resolved issue or expressed emotion, unable to move on.\\
$\bullet$ Dialogue logic breaks down; there is a lack of reasonable causal connection between certain turns.

\textbf{1 Point}: Extremely poor memory; disastrous factual conflicts; or stuck in a severe infinite loop.\\
\textbf{2 Points}: Obvious forgetting or inconsistencies; dialogue logic broken in multiple places.\\
\textbf{3 Points (Baseline)}: No obvious factual conflicts or forgetting. \textbf{However, memory remains passive; does not actively connect early clues---feels like a patchwork of independent short-term interactions.}\\
\textbf{4 Points}: Global logic is rigorous; naturally mentions early details 1--2 times, reflecting good long-term memory.\\
\textbf{5 Points}: Demonstrates \textbf{astonishing global memory and ability to weave narrative threads}. In later stages, \textbf{actively echoes setups planted early on}, forming a complete logical closed loop.

\textbf{Scoring Steps}\\
1. Read the full dialogue, mapping timeline and key details.\\
2. Start from 3 points.\\
3. Check deduction criteria: amnesia, contradictions, dead loops.\\
4. Look for bonus evidence of active callbacks and logical closure.\\
5. Write rationale, then provide score.

\textbf{Output Format}\\
\texttt{\{"context\_consistency": \{"reason": "<Around 150 words...>", "score": <Integer from 1-5>\}\}}
\end{tcolorbox}

\subsection{Prompt for User Persona Generation}
\label{app:prompt_user_persona}

\begin{tcolorbox}[
    breakable,
    colback=MyWordColor1,
    colframe=black,
    boxrule=0.6pt,
    arc=1.5mm,
    title={\textbf{Prompt for User Persona Generation}},
    fontupper=\small
]

\textbf{Case 1: Extracting the user persona from the role persona}

\medskip
\textbf{[System Prompt]}

You are a persona extraction assistant. Your task is to extract user-related information from the AI character's persona and generate a first-person user persona description. Note: you must generate the \textbf{user's} persona, not the AI character's persona.

\medskip
\textbf{[User Prompt]}

Below is the persona information of an AI character, which includes descriptions of the user who interacts with this character.

\medskip
Character Persona Information:
\{\textit{persona\_info}\}

\medskip
Task: Please carefully read the persona above, extract the relevant information about the \textbf{user/conversation partner}, and generate a first-person user persona description.

\medskip
Important reminders:
\begin{itemize}
    \item You are generating the \textbf{user's} persona, not the AI character's persona.
    \item If the persona says ``the user is XXX,'' then you should write ``I am XXX.''
    \item If the persona mentions the relationship between the user and the character, this relationship should be reflected in the generated persona.
    \item Use the first person (``I''); do not use the second person (``you'') or third person.
    \item Keep it within 80 words and make it concise and clear.
\end{itemize}

\bigskip
\hrule
\bigskip

\textbf{Case 2: Generating a generic user persona}

\medskip
\textbf{[System Prompt]}

You are a persona creation assistant. Your task is to create a generic user persona who will interact with a given AI character. Note: you are creating the \textbf{user}, not the AI character.

\medskip
\textbf{[User Prompt]}

Below is the persona information of an AI character:

\medskip
\{\textit{persona\_info}\}

\medskip
Task: Please generate a generic user persona for someone who will interact with the character above.

\medskip
Requirements:
\begin{itemize}
    \item Generate a \textbf{generic user persona}, not the AI character's persona.
    \item The user should know the character's basic background.
    \item The user should have their own personality traits (e.g., cheerful, introverted, humorous).
    \item Use the first person (``I'') and keep it within 60 words.
\end{itemize}

\end{tcolorbox}

\subsection{Prompt for User Simulator}

\begin{tcolorbox}[
    breakable,
    colback=MyWordColor1,
    colframe=black,
    boxrule=0.6pt,
    arc=1.5mm,
    title={\textbf{Prompt for User Simulator}},
    fontupper=\small
]

\textbf{[System Prompt]}

You are a user simulator. Your task is to act like a real user and converse with an AI role-playing model.
\textbf{Always remember: you are only the user. You must never play the other character's role.}

\medskip
\textbf{Key behavioral guidelines:}
\begin{enumerate}
    \item Behave like a lazy user: only passively answer the other party's questions, or give short comments on what they say. \textbf{Do not} proactively start new topics or ask questions frequently.
    \item \textbf{Colloquial style}: speak naturally and casually, like a real person texting.
    \item \textbf{Keep replies short}: each response should be strictly limited to about \textbf{10--20 words}, usually in \textbf{one sentence}.
\end{enumerate}

\medskip
\textbf{Your Persona (User Persona)}\\
Fully immerse yourself in the following persona:

[\{\textit{user\_system\_prompt}\}]

\medskip
\textbf{The AI you are talking to has the following persona}\\
(Note: this is \textbf{their} persona, \textbf{not yours}.)

[\{\textit{processed}\}]

\medskip
\textbf{Current Task}

Based on \textbf{your persona} and the \textbf{dialogue context}, give a short, natural, and passive reply.

\end{tcolorbox}

\twocolumn
